%------------------------------------------------
% Latex-Grundger�st f�r Seminarausarbeitungen
%
% zu Erzeugen mit
%   pdflatex template.tex
%
% von Carsten Gutwenger
% angepasst von Igor Vatolkin
%------------------------------------------------

\documentclass[a4paper,12pt,twoside]{article}
\usepackage[a4paper,left=3.5cm,right=2.5cm,bottom=3.5cm,top=3cm]{geometry}

% Encoding der .tex-Datei
% muss je nach Editor-Einstellungen eventuell angepasst werden!
%   latin1    Unix/Windows
%   ansinew   Alternative f�r Windows
%   applemac  Apple
%   utf8      Alternative f�r Unix/Windows/Apple
\usepackage[latin1]{inputenc}

%\usepackage[ngerman]{babel}                                 % deutsche Sprache
\usepackage[T1]{fontenc}                                    % Unterst�tzung f�r Umlaute mit Fonts
\usepackage[pdftex]{graphicx,color}                         % Einbetten von Grafiken, Farbe
\usepackage[format=plain,small,bf,indention=7.5mm]{caption} % verbesserte Beschriftung von Abbildungen
\usepackage{subfig}                                         % subfigures
\usepackage{url}                                            % \url{}-Kommando
\usepackage{fancyhdr,float}                                 % Kopf-/Fu�zeilen
\usepackage[pdftex,pdfpagelabels]{hyperref}                 % Hyperlinks im PDF
\usepackage{natbib,doi}                                     % Literaturverzeichnis, DOIs
\usepackage[section,boxed]{algorithm}                       % Floating-Umgebung f�r Algorithmen
\usepackage{algpseudocode}                                  % Pseudocode f�r Algorithmen

\captionsetup{farskip=10pt,topadjust=0pt,captionskip=10pt,nearskip=0pt,margin=10pt}

\pagestyle{fancy}
\fancyhead[LE,RO]{\thepage}
\fancyhead[RE]{\nouppercase{\slshape \leftmark}}
\fancyhead[LO]{\nouppercase{\slshape \rightmark}}
\fancyfoot[C]{}


\begin{document}

%------------------------------------------------
% Titelseite
%------------------------------------------------

\begin{titlepage}
\vspace*{-2cm}
\newlength{\links}
\setlength{\links}{-1.5cm} \sf \LARGE

\hspace*{\links}
\begin{minipage}{12.5cm}
\includegraphics[width=8cm]{RWTH_Logo_3.svg}
\end{minipage}

\vspace*{4cm}

\large
\begin{center}
{\Large Seminar work} \\[1ex]
{\LARGE\textbf{[Title]}}\\[3ex]
[Names]\\[1ex]
\today\\[7ex]
as part of seminar\\[1ex]
{\Large\textbf{Artificial Intelligence in Music Data Analysis}}\\[1ex]
organised by Dr.~Igor Vatolkin\\[1ex]
Summer term 2025
\end{center}

\vspace*{5cm}
\hspace*{\links}
\begin{minipage}[b]{15cm}
\normalsize \raggedright

\textbf{Based on:}\\
Authors, title of the paper, conference
\end{minipage}

\vfill
\hspace*{\links}
\begin{minipage}[b]{8cm}
\normalsize \raggedright
RWTH Aachen\\
Computer Science 14\\
Artificial Intelligence Methodology\\
\url{https://www.aim.rwth-aachen.de/}
\end{minipage}

\end{titlepage}

%------------------------------------------------
% Inhaltsverzeichnis
%------------------------------------------------

\tableofcontents
\clearpage


%------------------------------------------------
% Einf�hrung
%------------------------------------------------

\section{Introduction}

Lorem ipsum dolor sit amet, consetetur sadipscing elitr, sed diam nonumy eirmod tempor invidunt ut labore et dolore
magna aliquyam erat, sed diam voluptua. At vero eos et accusam et justo duo dolores et ea rebum. Stet clita kasd
gubergren, no sea takimata sanctus est Lorem ipsum dolor sit amet. Lorem ipsum dolor sit amet, consetetur sadipscing
elitr, sed diam nonumy eirmod tempor invidunt ut labore et dolore magna aliquyam erat, sed diam voluptua. At vero eos
et accusam et justo duo dolores et ea rebum. Stet clita kasd gubergren, no sea takimata sanctus est Lorem ipsum dolor
sit amet.

Example for a reference: \cite{HanKno06} (journal article) and \cite{TarekEtAl07} (conference paper).


%------------------------------------------------
% Literaturverzeichnis
%------------------------------------------------
\newpage
\bibliographystyle{plain}
\bibliography{literatur}
\addcontentsline{toc}{section}{\bibname}


\end{document}
